% !TEX program = xelatex
\documentclass{英语试卷}
\begin{document}
\examtitle{义乌市普通高中 2025 届第二次适应性考试}{英语}
\par 义乌市普通高中 2025 届第二次适应性考试\par
\par 英语试题\par
\begin{notice}
\par 1. 答卷前，考生务必将自己的姓名、准考证号填写在答题卡上。\par
\par 2. 回答选择题时，选出每小题答案后，用铅笔把答题卡上对应题目的答案标号涂黑。如需改动，用橡皮擦干净后，再选涂其他答案标号。回答非选择题时，将答案写在答题卡上，写在本试卷上无效。\par
\par 3. 考试结束后，将本试卷和答题卡一并交回。\par
\end{notice}
\parttitle{第一部分 听力（共两节，满分 30 分）}
\par 做题时，先将答案标在试卷上，录音内容结束后，你将有两分钟的时间将试卷上的答案转涂到答题卡上。微信公众号 IAI English\par
\sectitle{第一节（共5小题：每小题1.5分，满分7.5分）}
\par 听下面 5 段对话。每段对话后有一个小题，从题中所给的 A、B、C 三个选项中选出最佳选项。听完每段对话后，你都有 10 秒钟的时间来回答有关小题和阅读下一小题。每段对话仅读一遍。\par
\question{1. What did the man do yesterday evening?}
\begin{choices}(3)
  \item He played soccer.
  \item He saw a movie.
  \item He went shopping.
\end{choices}
\question{2. How does the man get to work?}
\begin{choices}(3)
  \item By bus.
  \item By subway.
  \item By car.
\end{choices}
\question{3. What's the probable relationship between the speakers?}
\begin{choices}(3)
  \item Workmates.
  \item Chef and waiter.
  \item Teacher and student.
\end{choices}
\question{4. Where will the man go first?}
\begin{choices}(3)
  \item The library.
  \item The baseball field.
  \item The supermarket.
\end{choices}
\question{5. How much should the man pay?}
\begin{choices}(3)
  \item £220.
  \item £196.
  \item £150.
\end{choices}
\sectitle{第二节（共15小题：每小题1.5分，满分22.5分）}
\par 听下面 5 段对话或独白。每段对话或独白后有几个小题，从题中所给的 A、B、C 三个选项中选出最佳选项。听每段对话或独白前，你将有时间阅读各个小题，每小题 5 秒钟；听完后，各小题将给出 5 秒钟的作答时间。每段对话或独白读两遍。\par
\par 听第 6 段材料，回答第 6、7 题。\par
\question{6. What are the speakers talking about?}
\begin{choices}(3)
  \item A car accident.
  \item A traffic report.
  \item A friend's experience.
\end{choices}
\question{7. What happened to Linda?}
\begin{choices}(3)
  \item She had a fall.
  \item She broke her arm.
  \item She was hit by a car.
\end{choices}
\par 听第 7 段材料，回答第 8、9 题。\par
\question{8. How does the woman feel now?}
\begin{choices}(3)
  \item Delighted and surprised.
  \item Excited and thankful.
  \item Happy and relieved.
\end{choices}
\question{9. Why does the woman want to visit Scotland?}
\begin{choices}(3)
  \item The scenery there is beautiful.
  \item She has a family member there.
  \item She wants to go to a warmer place.
\end{choices}
\par 听第 8 段材料，回答第 10 至 12 题。\par
\question{10. Why does Janet go to Detroit?}
\begin{choices}(3)
  \item To attend an exhibit.
  \item To go to a meeting.
  \item To visit some attractions.
\end{choices}
\question{11. What does Greg suggest Janet do?}
\begin{choices}(3)
  \item Join him on the train.
  \item Travel to Detroit by car.
  \item Take plenty of rest along the way.
\end{choices}
\question{12. When will the speakers leave for Detroit?}
\begin{choices}(3)
  \item On Monday.
  \item On Tuesday.
  \item On Saturday.
\end{choices}
\par 听第 9 段材料，回答第 13 至 16 题。\par
\question{13. What do we know about Mr. Theo?}
\begin{choices}(2)
  \item He used to be a famous actor. 3. He has contacts in the acting industry.
  \item He has remained friends with Joel Brookes.
\end{choices}
\question{14. What was Sasha's initial reaction to Joel Brookes' visit to her school?}
\begin{choices}(3)
  \item She was slightly doubtful.
  \item She was totally uninterested.
  \item She was pleasantly surprised.
\end{choices}
\question{15. Why did Joel Brookes come to Sasha's school?}
\begin{choices}(3)
  \item To meet the drama teacher.
  \item To choose shooting locations.
  \item To give a motivational speech
\end{choices}
\question{16. What was Sasha asked to do after meeting Mr. Theo?}
\begin{choices}(3)
  \item Practice a school play.
  \item Try out for a role in the film.
  \item Assist the producer in filming.
\end{choices}
\par 听第 10 段材料，回答第 17 至 20 题。\par
\question{17. How was the speaker's life after graduation?}
\begin{choices}(3)
  \item Creatively fulfilling.
  \item Emotionally rewarding.
  \item Financially challenging.
\end{choices}
\question{18. Who is Carly?}
\begin{choices}(3)
  \item A comedian.
  \item A fictional character.
  \item The speaker's childhood friend.
\end{choices}
\question{19. What shocked the speaker?}
\begin{choices}(3)
  \item The rejection for a job position.
  \item The unexpected success of the book.
  \item The approval from multiple publishers.
\end{choices}
\question{20. What inspired the speaker to create the book?}
\begin{choices}(3)
  \item Her mother's letters.
  \item Her grandmother's story.
  \item Her childhood experience.
\end{choices}
\parttitle{第二部分 阅读理解（共两节，满分 50 分）}
\sectitle{第一节（共15小题：每小题2.5分，满分37.5分）}
\par 阅读下列短文，从每题所给的 A、B、C 和 D 四个选项中，选出最佳答案，\par
\passagelabel{A}
\par Our diverse range of short online courses can be accessed through our partners FutureLearn and Coursera. Here's how our short online courses work.\par
\par Our short online courses\par
\par There are three collections of courses, categorized according to academic level:\par
\par 1. Online courses for high school and college students,\par
\par 2. Masters taster courses,\par
\par 3. Continuing professional development.\par
\par Our short courses include pre-recorded video lectures, readings and interactive activities. You can also interact and collaborate with fellow learners in discussions, forums and group activities.\par
\par Course materials can be accessed via web browsers and using the Coursera and Futurelearn mobile apps, compatible with smartphones and tablets Registering for a short online course\par
\par Once you have chosen your course, you will need to register through either the Coursera or FutureLearn website. Information on how to register can be found on the course page.\par
\par All of our online courses are available for free for a limited time, allowing you to explore different subjects and course materials without any financial commitment. The length of time you will have free access may differ between Coursera and FutureLearn, and the different courses. This information will be provided when you enroll.\par
\par If you wish to earn a certificate once you have completed your course, you can upgrade to obtain a course certificate, for a fee. Upgrading will also provide you with unlimited access to the course materials and assessments.\par
\question{21. What do we know about the short online courses?}
\begin{choices}(2)
  \item They are limited to masters only.
  \item They are grouped by academic stages.
  \item They are restricted to web browsers use.
  \item They are presented by live-streamed lectures.
\end{choices}
\question{22. What is the difference between registering on Coursera and FutureLearn?}
\begin{choices}(2)
  \item Course materials.
  \item Types of certificates.
  \item Payment methods.
  \item Free access duration.
\end{choices}
\question{23. What is the main purpose of the passage?}
\begin{choices}(2)
  \item To advertise paid certificate upgrades.
  \item To explain the benefits of online courses.
  \item To offer guidance on short online courses.
  \item To introduce two online course platforms.
\end{choices}
\passagelabel{B}
\par Italian cuisine may be one of the most recognized and loved in the world—often imagined as a collection of ancient recipes and local traditions preserved by skillful grandmothers over centuries. But now, a food historian is exposing common myths and false stories, arguing that the Italian cuisine as we know it is barely a few decades old.\par
\par Grandi claims in his latest book, Italian Cuisine Doesn't Exist, the idea that many beloved recipes and products such as cheeses or cured meats have hundreds of years of history is pure fantasy. Most dishes have a shorter history than most people imagine.\par
\par Migration is what made Italian cuisine what it is today. In the 19 th and 20 th Centuries, millions of people left Italy, a poverty-stricken country with limited dietary options, emigrating to South America, North America and European countries. The Italians who landed on New York's Ellis Island left hunger and misery behind. And it was in the "new world", Grandi argues, that Italian immigrants found the wealth and ingredients to create the recipes that eventually popularized Italian food worldwide.\par
\par However, many Italians are not familiar with this history. They might imagine that their beloved recipes were born in Italy, passed down, unchanged, from generation to generation and eventually exported abroad by Italian migrants.\par
\par Pizza is the most famous example of this. Born as a cheap street food in Naples, in the 19 Century it was synonymous with poverty and dirt. It was in America, says Grandi, that pizza became "red". While fresh tomatoes were among the original toppings, Italian immigrants to the US popularized pizza prepared with tomato sauce, a product that industrialization made easy to access and store. And it was in the US that pizzerias really started to take off. Scholars call this process "the pizza effect": when a product leaves its place of origin, gets profoundly transformed and then returns to its place of origin to be fully embraced in a completely different form.\par
\question{24. What does Grandi challenge in his book?}
\begin{choices}(2)
  \item The role of Italian chefs.
  \item The history of Italian cuisine.
  \item The recipe for cured meats.
  \item The process of cheese production.
\end{choices}
\question{25. How did Italian immigrants contribute to the spread of Italian cuisine?}
\begin{choices}(2)
  \item By running Italian restaurants worldwide.
  \item By promoting food industrialization in Italy.
  \item By creating new recipes with local ingredients.
  \item By preserving ancient recipes through generations.
\end{choices}
\question{26. What does the underlined phrase "was synonymous with" in the last paragraph mean?}
\begin{choices}(2)
  \item Equaled.
  \item Removed.
  \item Caused.
  \item Exposed.
\end{choices}
\question{27. Which of the following can be a suitable title for the text?}
\begin{choices}(2)
  \item The Migration of Italian Cuisine
  \item The Origin of the Pizza Effect
  \item The Secret of Traditional Recipes
  \item The Truths Behind Italian Cuisine
\end{choices}
\passagelabel{C}
\par Adriano Lippo is on his small boat, staring out over the dark blue waters of the Mar Piccolo. "I don't think it will happen this year," he says, shaking his head. "Everything is dead."\par
\par For the past decade, Lippo has made the same journey at the end of every summer, steering his boat to transfer mussels (贻贝) across the Mar Piccolo, a semi-enclosed bay located between the city and the open sea. Once the lifeblood of a centuries-old tradition, the Mar Piccolo supported generations of mussel farmers, producing 60,000 tons of mussels annually in the early 2000s. However, pollution and heatwaves have since taken a heavy toll, pushing the business into sharp decline and leaving farmers to pay the price.\par
\par Now, mussel farmers are pinning their hopes on a plan to restore the polluted Mar Piccolo through phytoremediation (植物修复). In early 2024, geologist Vito Uricchio and his team of researchers launched a pilot project, which focused on depolluting one hectare of polluted land using plants such as the Monviso clone. After a year of planting, chemical analyses revealed significant improvements, demonstrating the method's ability to address both organic and inorganic pollution. The process can treat soil up to 5m deep at a small part of the cost of traditional remediation methods.\par
\par Though current research so far shows promise, Werther Nissim, an assistant professor at the University of Milano-Bicocca, says phytoremediation alone is unlikely to fully depollute an area as large as the Mar Piccolo. He adds: "The need for a long time is another bottleneck. The technique should be combined with other more traditional remediation methods."\par
\par "Even with the best possible remediation, the benefits would be completely erased by the continuation of polluting activities," Uricchio remains full of fight. "I want to act; I want to remediate as much as possible."\par
\question{28. What has happened to mussel farming in the Mar Piccolo in recent years?}
\begin{choices}(2)
  \item It has suffered a serious drop
  \item It has lost the support of young people.
  \item It has failed to remain profitable.
  \item It has caused pollution in nearby waters.
\end{choices}
\question{29. What does paragraph 3 mainly talk about?}
\begin{choices}(2)
  \item The economic benefits of Monviso clones.
  \item Farmers' efforts to rebuild the Mar Piccolo.
  \item A test program using plants to clean polluted soil.
  \item Scientists' work on developing phytoremediation.
\end{choices}
\question{30. What is Nissim's attitude toward depolluting Mar Piccolo with phytoremediation?}
\begin{choices}(2)
  \item Cautious.
  \item Dismissive.
  \item Unclear.
  \item Appreciative.
\end{choices}
\question{31. What mainly ensures the success of remediation?}
\begin{choices}(2)
  \item Expanding area for Monviso clones.
  \item Reducing ongoing pollution actions.
  \item Mixing traditional and modern methods.
  \item Raising public environmental awareness.
\end{choices}
\passagelabel{D}
\par When questions about children arise, many parents turn to the internet and social media for advice on topics ranging from eating habits to children's screen time. But what kind of information are parents most likely to see and share?\par
\par A new study by NYU researchers, Erin O'Connor and Robin Neuhaus, examined how articles about screen time resonated with parents, and found that articles combining attention-grabbing language with credible scientific framing (表述) were the most likely to gain immediate and long-term attention on social media. O'Connor says: "Parents deserve information that not only captures their attention, but also supports informed decision-making about their children's development."\par
\par O'Connor and Neuhaus analyzed 136 articles published between 2016 and 2021. They used BuzzSumo, a social media analytics tool, to track each article's total social media shares across websites and its evergreen score—to find how much attention the article continued to receive after its initial publication.\par
\par Each article was evaluated for two main factors: sensationalism (煽动性) and scientific framing. Sensationalism was measured by the presence of warning language (e.g., "alarming," "skyrocket", "irreversible (不可逆转) damage"), while scientific framing was measured by how well the article explained findings, noted limitations, and gave accurate reporting.\par
\par Their analysis revealed that articles with high sensationalism and strong scientific framing achieved the highest immediate shares and long-term attention. In particular, articles that focused on younger children, with a negative view of screen time, and published on popular websites showed stronger performance. Interestingly, warning language alone boosted short-term sharing, but it was only when combined with strong scientific framing that articles kept people interested over time. Additionally, longer articles with more detailed context were more likely to achieve long-term attention.\par
\question{32. What does the study focus on?}
\begin{choices}(2)
  \item Children's use of mobile phones.
  \item Parents' reliance on online advice
  \item Methods to improve parenting skills.
  \item Features of popular screen-time articles.
\end{choices}
\question{33. How did researchers measure an article's lasting influence?}
\begin{choices}(2)
  \item By checking its shares in BuzzSumo.
  \item By calculating its evergreen score.
  \item By counting its total likes on websites.
  \item By predicting its long-term growth.
\end{choices}
\question{34. Which expression increases an article's sensationalism rating?}
\begin{choices}(2)
  \item Experts suggest limiting screen time for better sleep quality.
  \item Researchers noted that screen time effects vary by age group.
  \item Constant screen time produces non-repairable damage in kids.
  \item Parents can set an alarm to warn against too much screen time.
\end{choices}
\question{35. Which is a finding of the study?}
\begin{choices}(2)
  \item Negative content gets shared the most.
  \item Short articles attract more initial attention.
  \item Emotional language alone drives article popularity.
  \item Scientific framing influences articles' long-term popularity.
\end{choices}
\sectitle{第二节（共5小题；每小题2.5分，满分12.5分）}
\par 阅读下面短文，从短文后的选项中选出能填入空白处的最佳选项。选项中有两项为多余选项。 Medium friends are genuine; they share your history, circumstances or interests. They make you laugh, bring news, and offer insights. But, unlike your closest friends, they test the limits of your time, love, and energy. There are only so many dinners in a week, and so many people with whom you can be continuously texting. \blanknum{36}\par
\par And that is the problem with medium friends, the unclear boundaries we create with them. Without clear expectations, medium friendships can easily become unbalanced. You may like your medium friends more (or less) than they like you. \blanknum{37} But these conversations are harder with medium friends. 微信公众号 IAI English\par
\par \blanknum{38} They don't view it as a balance of give and take but as something that grows naturally, like art. Friendships, they say, are not ranked from best friends to strangers but reflect the effort two people put into them, showing what they enjoy, share, and how they're different.\par
\par Medium friends should not be seen as less valuable than best friends. \blanknum{39} They might be a tennis partner, a long-lost friend, or someone you play games with abroad. Without the pressure to be "perfect," these friendships can grow and serve each person as needed.\par
\par At their best, medium friendships can feel freeing, with less responsibility. Psychologist Beverley Fehr points to recent research showing that a wide, diverse friend group made up of weak and strong ties is best. She suggested, "\blanknum{40}" While the emotional and psychological benefits of medium friendships have not been studied, Fehr believes they are considerable.\par
\begin{choices}(1)
  \item Friendship is based on shared care and trust.
  \item Some social scientists see friendship differently.
  \item Don't put all your eggs in one relationship basket.
  \item They can be enjoyable and useful in their own way.
  \item With close friends, you can talk about issues or wounds.
  \item Scholars have tried to sort friendships into different levels.
  \item Medium friends provided the lie in trying to satisfy too many people.
\end{choices}
\parttitle{第三部分 语言知识运用（共两节，满分 30 分）}
\sectitle{第一节（共15小题；每小题1分，满分15分）}
\par 阅读下面短文，从短文后各题所给的 A、B、C 和 D 四个选项中，选出可以填入空白处的最佳选项。\par
\par Kitesurfers, riding the wind, appear to escape gravity. \blanknum{41} to a parasol (牵引伞) arching 40 feet high, they shift directions, slide along the waves and float through the air, seemingly \blanknum{42} and birdlike. The sport requires more than \blanknum{43} nerves. It also needs amazing athletic skills to \blanknum{44} both the wind and air at the same time-almost like superhuman. And for one man, his \blanknum{45} is also nearly superhuman. 微信公众号 IAI English\par
\par Bruno Lobo, a surgeon and a former champion swimmer, has been kitesurfing for 15 years. In 2016 and 2023, he held six Brazilian \blanknum{46} and won Olympic gold in kitesurfing.\par
\par On a \blanknum{47} day with some gray clouds and strong winds, Lobo went into the ocean near Sao Luis. "After \blanknum{48} a few meters, I heard a cry for help and noticed a girl who was drowning," he explained in a post on Instagram. "I \blanknum{49} approached her with the kite, tried to calm her down and asked her to \blanknum{50} on my back." Lobo used his equipment and experience as a \blanknum{51} and surfer to get the woman safely into the shallows, where lifeguards could \blanknum{52}.\par
\par CNN \blanknum{53} the story. But instead of shining the spotlight on himself, Lobo used the opportunity to educate people about the \blanknum{54} of watersports. Lobo reminded everyone that knowing our \blanknum{55} and adopting a little humility is just as crucial as the excitement of watersports.\par
\cz{41}{Glued}{Delivered}{Added}{Tied}
\cz{42}{fanciful}{dreamlike}{weightless}{peaceful}
\cz{43}{jumpy}{steely}{fragile}{tense}
\cz{44}{control}{inspect}{handle}{investigate}
\cz{45}{confidence}{toughness}{vision}{heroism}
\cz{46}{titles}{hospitals}{events}{positions}
\cz{47}{training}{racing}{working}{filming}
\cz{48}{swimming}{walking}{picturing}{riding}
\cz{49}{steadily}{immediately}{quietly}{simply}
\cz{50}{stamp}{play}{slide}{climb}
\cz{51}{lifeguard}{sailor}{swimmer}{surgeon}
\cz{52}{land}{assist}{respond}{arrive}
\cz{53}{picked up}{looked into}{made up}{settled on}
\cz{54}{laws}{tricks}{dangers}{impacts}
\cz{55}{hopes}{limits}{fears}{skills}
\sectitle{第二节（共10小题；每小题1.5分，满分15分）}
\par 阅读下面短文，在空白处填入 1 个适当的单词或括号内单词的正确形式，\par
\par Chinese mathematician Wang Hong, born in the 1990s, has resolved a century-old mathematical puzzle, the three-dimensional case of the Kakeya conjecture (猜想), establishing herself as a potential candidate \blanknum{56} the Fields Medal.\par
\par In 1917, Japanese mathematician Sōichi Kakeya posed a seemingly simple question: What is the \blanknum{57} (small) area that an infinitely (无限地) thin needle can sweep when rotated (旋转)? \blanknum{58} this problem was solved in two dimensions, its three-dimensional version remained unresolved for over a century, \blanknum{59} (puzzle) mathematicians worldwide\par
\par Wang, now 34, demonstrated \blanknum{60} (exception) talent from an early age. At just 16, she achieved outstanding results in China's national college entrance exam, earning \blanknum{61} (admit) to Department of Geosciences at the Peking University. \blanknum{62} (drive) by a deep interest in mathematics, she later transferred to the university's mathematics department.\par
\par During her doctoral studies at the Massachusetts Institute of Technology, Wang studied under renowned mathematician Larry Guth, \blanknum{63} leading authority in geometric measure theory, \blanknum{64} guidance greatly shaped her academic development. Since July 2023, Wang has served as an associate professor at New York University.\par
\par In February 2025, Wang, together with her collaborator Joshua Zahl, \blanknum{65} (release) a 127-page paper formally announcing their achievement, which has gained considerable attention from the global mathematical community.\par
\parttitle{第四部分 写作（共两节，满分 40 分）}
\sectitle{第一节（满分15分）}
\par 假定你是李华，你最近参加了一场关于“职业生涯规划”的讲座。请你给英国朋友 Michael 写一封邮件分享这次经历，内容包括：(1) 讲座的主要内容；(2) 你对职业规划的看法。\par
\par 注意：(1) 写作词数应为 80 左右；(2) 可适当增加细节，以使行文连贯。\par
\par Dear Michael,\par
\par I'm writing to share with you a career planning seminar I recently attended.\par
\sectitle{第二节（满分25分）}
\par 阅读下面材料，根据其内容和所给段落开头语续写两段，使之构成一篇完整的短文。\par
\par In the summer of 2000, I was barely fifteen years old. I wore glasses, was a bit overweight, and was incredibly insecure and shy. Most of my peers thought I wasn't worth much and could never accomplish anything.\par
\par By that spring, I had withdrawn almost completely and preferred to be by myself, reading or drawing. My parents were worried and offered me the chance to attend a three-week summer camp in Orr, Minnesota.\par
\par Eager to escape, I jumped at the opportunity.\par
\par The campsite on the Pelican Lake had plenty to offer: acres of forest, bunkhouses, a mess hall, a gymnasium, shower facilities, and a full schedule of activities. But one experience stood out.\par
\par "Hey, Jojo! Ever go ‘Banana-Rafting (漂流)’?" Greg, our counselor (辅导员), looked around at our group of fifteen campers and settled on me.\par
\par "No, what's that?"\par
\par "Here, I’ll demonstrate!" He held up two bright yellow lifejackets. He put on one lifejacket normally and then the other upside down. He stepped his left leg through one armhole and right leg through the other, straddling the neck. He pulled the lifejacket up, zipped it up, and belted it around his hips. Turning around like a model, he showed us all what he looked like: a waddling yellow banana with arms, legs, and a head! 微信公众号 IAI English\par
\par "Watch!" He took a running jump off an embankment (堤坝) into the riffles (浅滩) before the rapids on the river that fed Pelican Lake. Horrified and excited, we rushed to the edge to watch as he floated down, grabbed a rope, and pulled himself to the shore.\par
\par "That's Banana-Rafting!" Greg said, dripping wet but happy. All of us yelled and stormed the pile of yellow lifejackets.\par
\par "Are you ready?" Greg asked as I nervously stood on the edge overlooking the river.\par
\par "No—No I'm not!"\par
\par "It's okay to be scared," Greg continued. "You're brave and strong. We all believe in you."\par
\par 注意:\par
\par (1) 续写词数应为 150 左右;\par
\par (2) 请按如下格式在答题卡的相应位置作答。\par
\writingpara{Paragraph 1: Everyone behind me was cheering me on.}
\writinglines{8}
\writingpara{Paragraph 2: Something changed in me that day. \blanknum{}}
\writinglines{8}
\end{document}
